\documentclass[reqno]{exam}

\usepackage{amssymb, amsmath, amsthm, stmaryrd}
\usepackage{fullpage, parskip}
\usepackage{enumitem}
\usepackage{tikz}
\usepackage{ulem}
\usepackage{hyperref}

\title{Introduction to Computational Math Spring 2026 \\ Homework 05}
\date{Due: Friday, Feb 27, 2026, 11:59 PM}
\author{}

\begin{document}
\maketitle
\thispagestyle{empty}

Textbook = Driscoll and Braun (\url{https://fncbook.com/})

There are more problems in this homework than the previous one, but they are shorter and more conceptual.

\begin{enumerate}


\item Textbook Exercise 3.2.1

\item Textbook Exercise 3.2.2

\item Textbook Exercise 3.2.3

\item Textbook Exercise 3.2.5a

\item Textbook Exercise 3.4.1

\item Textbook Exercise 3.4.2

\item Textbook Exercise 3.4.3

To show that a square matrix $Q$ is orthogonal, it is enough to show that $Q^T Q = I$.

\item Textbook Exercise 3.4.7a

This question is asking you to find $\theta$ in terms of $\alpha$ and $\beta$. Your answer will involve a trigonometric function. 

\end{enumerate}

No Jupyter notebook this week.



\section*{Quiz Information}

The quiz will be held on {\bf Tuesday, March 03, 2026} during the discussion section.

Quiz questions will be modifications of the homework problems and material discussed in class.

\textbf{Topics covered:} 

Chapters 3.2 and 3.4. 

You will need memorize 
\begin{itemize}
    \item The normal equations for least squares problems,
    \item The formula for the Householder reflector.
\end{itemize}

You should expect both computational and conceptual questions on the quiz.

There won't be any questions about using the normal equations in terms of $QR$ factorization, but you should know how to use the normal equations in terms of the $A^T A$ matrix.



\end{document}